% Chapter5/chapter5.tex -- Results and Discussion
% -----------------------------------------------------------------------------
% This is the chapter your examiner reads most carefully. Every claim must be
% backed by a number from an experiment you actually ran.
%
% Suggested structure:
%   5.1 -- Experimental setup (so results are reproducible)
%   5.2 -- HLS synthesis results (resource utilisation, timing)
%   5.3 -- Functional verification
%   5.4 -- Performance benchmarks (throughput, latency)
%   5.5 -- Energy efficiency
%   5.6 -- Comparison against baselines
%   5.7 -- Discussion (interpret your results -- don't just describe them)

\chapter{Results and Discussion}
\label{ch:results}

\section{Experimental Setup}
\label{sec:exp_setup}

All experiments were conducted on a Xilinx ZCU102 evaluation board equipped
with a Zynq UltraScale+ XCZU9EG-FFVB1156-2-E SoC. The programmable logic
was clocked at \SI{200}{\mega\hertz} and the PS ran Ubuntu 20.04 on the
ARM Cortex-A53 at \SI{1.2}{\giga\hertz}. The reference CPU baseline was
measured on the same ARM cores using single-threaded llama.cpp with INT8
weights to ensure a fair comparison. Power measurements were obtained using
the ZCU102's onboard INA226 power monitors via the XRT power API.

\section{HLS Synthesis Results}
\label{sec:synthesis_results}

\Cref{tab:resource_util} summarises the FPGA resource utilisation reported
by Vitis HLS after synthesis and implementation in Vivado. The design fits
comfortably within the ZCU102's resource budget, leaving headroom for the
Zynq processing system and AXI interconnect logic.

\begin{table}[htbp]
  \centering
  \caption{FPGA resource utilisation after Vivado implementation on the ZCU102
           (xczu9eg-ffvb1156-2-e). Percentages are relative to total available
           resources on the device.}
  \label{tab:resource_util}
  \begin{tabular}{lrrrr}
    \toprule
    Resource & Used & Available & Utilisation (\%) \\
    \midrule
    LUT      & -- & 274,080 & -- \\
    FF       & -- & 548,160 & -- \\
    BRAM     & -- &     912 & -- \\
    DSP      & -- &   2,520 & -- \\
    \bottomrule
  \end{tabular}
\end{table}

% -- Note to Aaron: fill in the "Used" column once you have synthesis results.
% The "--" placeholders will compile cleanly in the meantime.

\lipsum[1] % Replace with your analysis of the synthesis report

\section{Functional Verification}
\label{sec:functional_verification}

\lipsum[2] % Describe your C simulation results, cosim, and on-board verification

\section{Inference Performance}
\label{sec:performance}

\begin{table}[htbp]
  \centering
  \caption{Inference performance comparison between the FPGA implementation
           and a CPU software baseline. Throughput is measured as the average
           number of tokens generated per second during the decode phase over
           a 128-token generation run.}
  \label{tab:performance}
  \begin{tabular}{lrrrr}
    \toprule
    Platform & Throughput (tok/s) & Latency/tok (ms) & Power (W) & Efficiency (tok/J) \\
    \midrule
    ARM Cortex-A53 (SW) & -- & -- & -- & -- \\
    ZCU102 FPGA (this work) & -- & -- & -- & -- \\
    \bottomrule
  \end{tabular}
\end{table}

\lipsum[3] % Replace with your performance analysis

\section{Energy Efficiency}
\label{sec:energy}

\lipsum[4] % Replace with your power and energy efficiency analysis

\section{Comparison Against Prior Work}
\label{sec:comparison}

\begin{table}[htbp]
  \centering
  \caption{Comparison of this work against prior FPGA-based transformer
           accelerators reported in the literature.}
  \label{tab:comparison}
  \begin{tabular}{llrrrr}
    \toprule
    Work & Platform & Model & Throughput & Power & Precision \\
    \midrule
    \cite{PLACEHOLDER} & -- & BERT  & -- & -- & INT8 \\
    \cite{PLACEHOLDER} & -- & GPT-2 & -- & -- & FP16 \\
    This work & ZCU102 & TinyLlama-1.1B & -- & -- & INT8 \\
    \bottomrule
  \end{tabular}
\end{table}

\lipsum[5] % Replace with your comparison analysis

\section{Discussion}
\label{sec:discussion}

\lipsum[6] % Replace with your interpretation of results, trade-offs, limitations
% Good discussion covers: why did you get these results? what surprised you?
% what are the bottlenecks? what would you do differently?
